\section{Working with Keyword Arguments}
Keyword arguments solve a common readability problem. When a function accepts many parameters, a long positional call can become difficult to read. Unlike positional arguments, keyword arguments allow you to specify exactly which parameter is being assigned a value, enhancing readability and reducing errors. A keyword call makes the meaning visible because the parameter names appear directly in the function call.

In this section, students will keep the same classifier behavior but improve the way the program is controlled from outside the function. This is where the implementation tasks begin.


\subsubsection{Advantages of Keyword Arguments:}

There are several advantages:
\begin{itemize}\setlength{\itemsep}{0pt}
    \item Enhanced readability and clarity in function calls.
    \item Flexible argument ordering, reducing potential errors.
    \item Easy handling of optional parameters with default values.

\subsubsection{Keyword Argument Syntax}
When you use keyword arguments, you specify the name of the parameter followed by an equal sign and the value it should receive within the function call. For example,
\begin{center}
    width=1000
\end{center}

\noindent
The benefit of this approach is that you can rearrange the order of arguments when calling the function without affecting its behavior, as long as all required arguments are provided.

So far, we have not been using keyword arguments in our code, but we will start introducing them in this session to give you more control over the behavior of your functions without having to change their internal logic.


\subsection{Task 1: Refactor \texttt{calculate\_accuracy} to use keyword-friendly defaults}

At first, this may seem unnecessary because \texttt{calculate\_accuracy} only uses two parameters. However, this is still a very useful practice for beginners. The goal of this task is not only to make this one function work, but also to help you understand how keyword arguments behave in Python.

When a function uses keyword-friendly parameters:
\begin{itemize}
    \item the function call becomes easier to read,
    \item the meaning of each value becomes clearer,
    \item and the function becomes safer to use later when the program grows larger.
\end{itemize}

In this task, you will make a small refactor:
\begin{itemize}
    \item you will update the function header so both parameters have default values,
    \item and then you will call the function using keyword arguments instead of a positional-style call.
\end{itemize}

Do not change the main accuracy logic. The purpose here is to improve the function interface, not to redesign how the calculation works.

\begin{tcolorbox}[title=Task 1: Refactor \texttt{calculate\_accuracy} to use keyword-friendly defaults]
Work only in \texttt{session5/session5\_iris\_template.py}. Do not create a new file, and do not rename the function.

\begin{enumerate}
    \item Find the function named \texttt{calculate\_accuracy}. Read the function carefully before editing anything so you understand what it already does.
    \item Locate the \texttt{<your\_code>} marker in the function header.
    \item Replace only the function header part so it becomes:
    \texttt{def calculate\_accuracy(correct=0, total=0):}
    \item Notice what this change means:
    \begin{itemize}
        \item \texttt{correct=0} gives the parameter \texttt{correct} a default value of \texttt{0},
        \item \texttt{total=0} gives the parameter \texttt{total} a default value of \texttt{0},
        \item so Python can still call the function even if no argument is passed in.
    \end{itemize}
    \item Keep the existing calculation logic and return statement inside the function body. Do not remove the accuracy logic that is already there.
    \item After updating the function definition, go to \texttt{setup\_application(...)}.
    \item Find the place where \texttt{calculate\_accuracy(...)} is called.
    \item Change that call so it uses keyword arguments:
    \texttt{correct=metrics["correct"]} and \texttt{total=metrics["total"]}.
    \item Read the new function call slowly:
    \begin{itemize}
        \item the value from \texttt{metrics["correct"]} is being assigned to the parameter named \texttt{correct},
        \item the value from \texttt{metrics["total"]} is being assigned to the parameter named \texttt{total}.
    \end{itemize}
    \item Quick check: \texttt{calculate\_accuracy(correct=8, total=10)} should return \texttt{80.0}.
    \item Quick check: \texttt{calculate\_accuracy()} should not crash. Because the function now has default values, Python can call it even when no arguments are given.
\end{enumerate}
\end{tcolorbox}

Use this exact function shape as the target:

\begin{minted}[fontsize=\small, linenos, breaklines=true, frame=lines]{python}
def calculate_accuracy(correct=0, total=0):
    # <your_code>: keep the accuracy logic and return value here
\end{minted}

\begin{tcolorbox}[title=First Unit Test for Task 1]
If this is your first time using \texttt{unittest}, think of it as an automatic checker for your Python code. Instead of checking your function by hand every time, Python can run several checks for you and report whether your function behaves correctly.

For Task 1, use the dedicated test file:
\texttt{session5/test\_calculate\_accuracy.py}

This test file checks only \texttt{calculate\_accuracy}. That is useful because you may still be working on the rest of \texttt{session5/session5\_iris\_template.py}.

Run the test from the project root folder with this command:

\begin{minted}[fontsize=\small, linenos, breaklines=true, frame=lines]{bash}
python3 -m unittest session5.test_calculate_accuracy -v
\end{minted}

What each part means:
\begin{itemize}
    \item \texttt{python3} starts Python,
    \item \texttt{-m unittest} runs Python's built-in testing tool,
    \item \texttt{session5.test\_calculate\_accuracy} points to the test file,
    \item \texttt{-v} means ``verbose'' and shows each test name clearly.
\end{itemize}

How to read the result:
\begin{itemize}
    \item \texttt{OK} means all tests passed,
    \item \texttt{FAIL} means the function ran, but at least one answer was wrong,
    \item \texttt{ERROR} usually means Python could not read or run the function yet.
\end{itemize}

For this task, the most important checks are:
\begin{itemize}
    \item \texttt{calculate\_accuracy(correct=8, total=10)} should return \texttt{80.0},
    \item \texttt{calculate\_accuracy()} should return \texttt{0.0} without crashing,
    \item keyword calls and positional calls should match when the numbers are the same.
\end{itemize}
\end{tcolorbox}

This code means:
\begin{itemize}
    \item the function is still named \texttt{calculate\_accuracy},
    \item the parameter \texttt{correct} defaults to \texttt{0},
    \item the parameter \texttt{total} defaults to \texttt{0},
    \item and the existing logic inside the function should stay responsible for computing and returning the result.
\end{itemize}

Then, inside \texttt{setup\_application(...)}, update the function call to use keyword arguments:

\begin{minted}[fontsize=\small, linenos, breaklines=true, frame=lines]{python}
accuracy = calculate_accuracy(
    correct=metrics["correct"],
    total=metrics["total"],
)
\end{minted}

Study this call carefully:
\begin{itemize}
    \item \texttt{correct=} tells Python exactly which parameter should receive \texttt{metrics["correct"]},
    \item \texttt{total=} tells Python exactly which parameter should receive \texttt{metrics["total"]},
    \item so the call is easier to read than simply writing two bare values in order.
\end{itemize}

The main advantage of this approach is that the parameter names are visible at the call site. A reader does not need to guess what each value means. Another advantage is that, when using keyword arguments, the order can be changed without changing the behavior, as long as the correct parameter names are used.

For example, you can call \texttt{calculate\_accuracy(total=total, correct=correct)} and it will work just fine.




\begin{tcolorbox}[title=Line-by-Line Debugging for Task 1]
Set a breakpoint on the line that calls \texttt{calculate\_accuracy(...)}.
\begin{itemize}
    \item Inspect \texttt{metrics["correct"]} and \texttt{metrics["total"]}.
    \item Use \textbf{Step Into} and confirm that those values arrive in the parameters \texttt{correct} and \texttt{total}.
    \item Tell students to notice that keyword argument order can change without changing the result.
\end{itemize}
This is one of the simplest examples for teaching the difference between positional matching and keyword matching.
\end{tcolorbox}




\subsection{Task 2: Build the keyword-friendly \texttt{setup\_application} function}
It is also possible to have many keyword arguments in a function.
For example, in the \texttt{setup\_application} function, you can have multiple keyword arguments to control different aspects of ... 

Recall that in previous session the setup application function accept 6 positional arguments, which means the order of the arguments matters. 

\begin{tcolorbox}[title=Task 2: Build the keyword-friendly \texttt{setup\_application} function]
In the template, find \texttt{setup\_application}.

Replace the \texttt{<your\_code>} marker in the function header so all four parameters become optional and keyword-friendly:
\texttt{def setup\_application(filepath=None, threshold=None, print\_each=True, use\_fit=False):}
\end{tcolorbox}

Use this exact function shape as the target:

\begin{minted}[fontsize=\small, linenos, breaklines=true, frame=lines]{python}
def setup_application(filepath=None, threshold=None, print_each=True, use_fit=False):
    # <your_code>: load the dataset, read the default settings, and apply the threshold rules
\end{minted}

Inside the function, the threshold logic must follow this priority order:
\begin{enumerate}
    \item if \texttt{use\_fit} is \texttt{True}, learn the threshold from \texttt{fit\_threshold\_from\_setosa(...)},
    \item else if \texttt{threshold} is not \texttt{None}, use the manual threshold,
    \item else keep the default threshold from \texttt{get\_default\_settings()}.
\end{enumerate}

If \texttt{filepath} is omitted, use \texttt{session5/iris\_data.csv} so the default run works without extra arguments.

Put this logic inside the function:

\begin{minted}[fontsize=\small, linenos, breaklines=true, frame=lines]{python}
dataset = load_iris_data(filepath)
settings = get_default_settings()

# <your_code>: apply the threshold priority order here

metrics = initialize_predictions(dataset, settings, print_each)
accuracy = calculate_accuracy(
    correct=metrics["correct"],
    total=metrics["total"],
)
return settings, dataset, metrics, accuracy
\end{minted}

This function is the main refactoring target of Session 5. It turns several separate steps into one reusable entry point.

\begin{tcolorbox}[title=Checkpoint]
After Task 2, students should be able to explain the difference between these two cases:
\begin{itemize}
    \item \texttt{threshold=2.2} means ``use this manual threshold'',
    \item \texttt{use\_fit=True} means ``ignore the manual override and learn the threshold from the dataset''.
\end{itemize}
\end{tcolorbox}

\begin{tcolorbox}[title=Line-by-Line Debugging for Task 2]
This is the most important debugging task in Session 5.
\begin{itemize}
    \item Set a breakpoint inside \texttt{setup\_application(...)} before the threshold logic.
    \item Inspect \texttt{threshold}, \texttt{use\_fit}, and \texttt{settings["threshold"]}.
    \item Step line by line through the \texttt{if use\_fit} and \texttt{elif threshold is not None} branches.
    \item Ask students: ``Which branch ran, and why?''
\end{itemize}
This question forces students to connect the function call in \texttt{main()} with the branch behavior inside the function.
\end{tcolorbox}

\subsection{Task 3: Add the default run in \texttt{main()}}
\begin{tcolorbox}[title=Task 3: Add the default run in \texttt{main()}]
In the template, find the default-run block and replace the \texttt{<your\_code>} marker by uncommenting the call.

The run should use the simplest possible call:
\texttt{setup\_application()}
\end{tcolorbox}

\begin{minted}[fontsize=\small, linenos, breaklines=true, frame=lines]{python}
# <your_code>: uncomment the default run block below
# make_print_status("Default run")
# settings, dataset, metrics, accuracy = setup_application()
# print_summary("Default run", settings, metrics, accuracy)
\end{minted}

This run is important because it proves that the defaults are working together:
\begin{itemize}
    \item default filepath,
    \item default threshold,
    \item default \texttt{print\_each=True},
    \item default \texttt{use\_fit=False}.
\end{itemize}

\begin{tcolorbox}[title=Checkpoint]
If the default run fails, students should first check the function headers. In beginner code, many default-run bugs come from forgetting to add \texttt{=None}, \texttt{=True}, or \texttt{=False} in the function definition.
\end{tcolorbox}

\subsection{Task 4: Add the positional override run in \texttt{main()}}
\begin{tcolorbox}[title=Task 4: Add the positional override run in \texttt{main()}]
In the template, find the positional run block and replace the \texttt{<your\_code>} placeholder by uncommenting the call.

Students should use a positional call like this:
\end{tcolorbox}

\begin{minted}[fontsize=\small, linenos, breaklines=true, frame=lines]{python}
# <your_code>: uncomment the positional override run below
# make_print_status("Positional override run")
# settings, dataset, metrics, accuracy = setup_application(
#     "session5/iris_data.csv", 1.8, False, False
# )
# print_summary("Positional override run", settings, metrics, accuracy)
\end{minted}

Explain to students that the order of positional arguments matters:
\begin{itemize}
    \item first argument $\rightarrow$ \texttt{filepath}
    \item second argument $\rightarrow$ \texttt{threshold}
    \item third argument $\rightarrow$ \texttt{print\_each}
    \item fourth argument $\rightarrow$ \texttt{use\_fit}
\end{itemize}

\begin{tcolorbox}[title=Checkpoint]
If students accidentally swap \texttt{1.8} and \texttt{False}, the run will still be syntactically valid, but the behavior will be wrong. This is exactly why positional calls become harder to read as the number of arguments grows.
\end{tcolorbox}

\begin{tcolorbox}[title=Line-by-Line Debugging for Task 4]
Set a breakpoint on the positional call in \texttt{main()} and use \textbf{Step Into}.
\begin{itemize}
    \item Before stepping in, read the four argument values left to right.
    \item After stepping in, compare them with the parameter names in \texttt{setup\_application(...)}.
    \item Inspect \texttt{filepath}, \texttt{threshold}, \texttt{print\_each}, and \texttt{use\_fit}.
\end{itemize}
This side-by-side comparison is one of the best debugging exercises in the whole session.
\end{tcolorbox}

\subsection{Task 5: Add the keyword override run in \texttt{main()}}
\begin{tcolorbox}[title=Task 5: Add the keyword override run in \texttt{main()}]
In the template, find the keyword-run block and replace the \texttt{<your\_code>} marker by uncommenting the call.

Students should use a named call like this:
\end{tcolorbox}

\begin{minted}[fontsize=\small, linenos, breaklines=true, frame=lines]{python}
# <your_code>: uncomment the keyword override run below
# make_print_status("Keyword override run")
# settings, dataset, metrics, accuracy = setup_application(
#     threshold=2.2,
#     print_each=False,
# )
# print_summary("Keyword override run", settings, metrics, accuracy)
\end{minted}

Explain to students that the order of keyword arguments does not matter as long as the parameter names are correct. For example, the following is also valid:

\begin{minted}[fontsize=\small, linenos, breaklines=true, frame=lines]{python}
setup_application(
    print_each=False,
    threshold=2.2,
)
\end{minted}

\begin{tcolorbox}[title=Checkpoint]
Students should notice that the keyword override run is easier to read than the positional override run. The call itself explains the meaning of each changed value.
\end{tcolorbox}

\subsection{Task 6: Add the fit-based run in \texttt{main()}}
\begin{tcolorbox}[title=Task 6: Add the fit-based run in \texttt{main()}]
In the template, find the fit-based block and replace the \texttt{<your\_code>} marker by uncommenting the call.

Students should use:
\end{tcolorbox}

\begin{minted}[fontsize=\small, linenos, breaklines=true, frame=lines]{python}
# <your_code>: uncomment the fit-based run below
# make_print_status("Fit-based run")
# settings, dataset, metrics, accuracy = setup_application(
#     print_each=False,
#     use_fit=True,
# )
# print_summary("Fit-based run", settings, metrics, accuracy)
\end{minted}

This run introduces a very important programming idea: one flag can change the path of the program. The data loading stays the same. The prediction loop stays the same. Only the threshold selection behavior changes.

\begin{tcolorbox}[title=Checkpoint]
After Task 6, students should compare:
\begin{itemize}
    \item the default run threshold,
    \item the positional override threshold,
    \item the keyword override threshold,
    \item and the fit-based threshold.
\end{itemize}
This comparison helps students see that Session 5 is about controlling one pipeline in multiple ways.
\end{tcolorbox}

\begin{tcolorbox}[title=Common Argument Mistakes]
Students often hit the same beginner errors in this session:
\begin{itemize}
    \item \textbf{Missing required positional argument}: the function header was not updated with a default value.
    \item \textbf{SyntaxError: positional argument follows keyword argument}: a positional value was placed after a keyword argument in the same call.
    \item \textbf{TypeError: multiple values for argument}: the same parameter was filled once positionally and once by keyword.
    \item \textbf{Wrong branch in \texttt{setup\_application(...)}:} students expected \texttt{threshold=} to run, but \texttt{use\_fit=True} took priority.
\end{itemize}
\end{tcolorbox}
